\documentclass[10pt,twocolumn]{article}

\usepackage[T1]{fontenc}
\usepackage[utf8]{inputenc}
\usepackage{lmodern}
\usepackage{microtype}
\usepackage[a4paper,margin=0.68in,columnsep=0.25in]{geometry}
\usepackage{amsmath,amssymb}
\usepackage{booktabs}
\usepackage{tabularx}
\usepackage{array}
\usepackage{enumitem}
\usepackage{xcolor}
\usepackage{url}
\usepackage{hyperref}
\usepackage{balance}

\urlstyle{same}
\hypersetup{
  colorlinks=true,
  linkcolor=blue!45!black,
  citecolor=blue!45!black,
  urlcolor=blue!55!black,
  pdftitle={Preference Post-Training Qwen2.5-0.5B},
  pdfauthor={Dheeraj Dhillon}
}

\setlength{\parindent}{1em}
\setlength{\parskip}{0pt}
\setlist{nosep,leftmargin=1.3em}
\newcommand{\model}{\textsc{Qwen2.5-0.5B-Instruct}}
\newcommand{\dataset}{\textsc{HelpSteer3-Preference}}
\newcommand{\trl}{\textsc{TRL}}
\newcolumntype{Y}{>{\raggedright\arraybackslash}X}

\title{\vspace{-1.5cm}\textbf{Preference Post-Training Qwen2.5-0.5B}\\
\large Supervised Fine-Tuning, Reward Modeling, PPO, and Auditable Evaluation}
\author{Dheeraj Dhillon}
\date{July 2026}

\begin{document}
\maketitle

\begin{abstract}
This work studies a complete reinforcement learning from human feedback
(RLHF) pipeline on \model{} with NVIDIA's \dataset{}. Although the starting
model is already instruction-tuned, one epoch of response-only supervised
fine-tuning (SFT) is used as domain and preference adaptation and as the common
reference policy for later alignment. A two-epoch scalar reward model and a
KL-regularized Proximal Policy Optimization (PPO) policy are then trained with
long-context, EOS-safe preprocessing and implementation details drawn from the
N+ RLHF reproduction.

SFT reaches 72.02\% teacher-forced completion-token accuracy on 1,917 retained
validation pairs. The reward model ranks the human-preferred response correctly
in 65.62\% of those pairs. On a separate generation suite covering all 2,017
validation prompts, the learned reward model prefers PPO over the original
instruction model in 1,027 cases, prefers Base in 981, and assigns 9 ties
(50.92\% PPO win rate). This narrow proxy-judge gain is accompanied by
substantially greater length, cap hits, and repetition. The result is therefore
presented as an inspectable alignment case study, not as evidence that PPO
universally improves the base model. Complete outputs and a balanced curated
subset are published for direct review.
\end{abstract}

\section{Introduction}

RLHF attempts to connect next-token language modeling with human judgments of
helpfulness, relevance, correctness, clarity, and safety. Its canonical
pipeline first adapts a policy to preferred demonstrations, learns a scalar
reward from response comparisons, and then optimizes the policy under a
constraint to its supervised reference \cite{stiennon2020,ouyang2022}.

This project asks whether that pipeline can be implemented, measured, and
audited on a half-billion-parameter model using rented notebook compute. The
small model makes iteration accessible, but it also exposes capacity limits
and makes reward-model error consequential. The Qwen2.5 model family is
described in \cite{qwen25}. The work began with custom
SFT, reward-model, and token-level PPO loops, then migrated the trainer
mechanics to Hugging Face \trl{} while retaining repository-owned data
contracts, model initialization, checkpoint records, evaluation, and
qualitative audit.

\paragraph{Why SFT an instruction model?}
The starting checkpoint is already instruction-tuned, so this stage is not
instruction tuning from scratch. It is continued SFT on the preferred
\dataset{} responses. This is justified for two reasons: it adapts the
response distribution to the exact preference data used downstream, and it
creates a reproducible reference policy shared by reward modeling, PPO, and
the planned DPO comparison. Using the original instruction checkpoint directly
would remove this controlled reference and make downstream differences harder
to attribute. Continued supervised adaptation before preference optimization
is also consistent with modern post-training pipelines \cite{ouyang2022}.

The principal contribution is not a state-of-the-art model. It is an
end-to-end study in which training and failure analysis receive equal weight:
full responses, reward scores, length, EOS behavior, cap hits, repeated
4-grams, deterministic triage, and curated examples remain available for
inspection.

\section{Method}

\subsection{Data and preprocessing}

\dataset{} contains 40,476 pairwise records over general, code, STEM, and
multilingual tasks \cite{wang2025}. Each record contains one conversational
context, two candidate responses, an overall preference in
$\{-3,\ldots,3\}$, and up to three individual annotator judgments. Negative
overall values favor response 1, positive values favor response 2, and zero is
a tie.

Tied or invalid pairs are excluded from preference training, leaving 36,264
training and 1,917 validation records. All stages use the Qwen chat template
and distinct PAD and EOS tokens. Completions are preserved through EOS; when a
sequence exceeds 4,096 tokens, old prompt turns are removed before a final
prompt-side token fallback. This reduces observed SFT training truncation from
38.47\% at a 1,024-token budget to 0.83\% at 4,096.

\subsection{Supervised fine-tuning}

For prompt $x$, preferred response $y=(y_1,\ldots,y_T)$, and completion mask
$m_t$, response-only SFT minimizes
\begin{equation}
\mathcal{L}_{\mathrm{SFT}}(\theta)
=-\sum_{t=1}^{T}m_t\log\pi_\theta(y_t\mid x,y_{<t}).
\label{eq:sft}
\end{equation}
Prompt and padding tokens do not contribute to the loss. LoRA
\cite{hu2022} rank-16 adapters are trained in all attention and feed-forward
projections; the backbone remains frozen. The merged result supplies both the
initial policy and the fixed reference for preference optimization.

\subsection{Reward modeling}

The reward model begins from the merged SFT backbone and adds a scalar head.
For preferred $y^+$ and rejected $y^-$ under the same prompt, it minimizes the
Bradley--Terry loss
\begin{equation}
\mathcal{L}_{\mathrm{RM}}(\phi)
=-\log\sigma\!\left(r_\phi(x,y^+)-r_\phi(x,y^-)\right).
\label{eq:rm}
\end{equation}
Only score differences are identifiable; the sign of one isolated reward is
not a quality boundary. A fixed offset estimated on preferred demonstrations
is therefore subtracted consistently during PPO and evaluation.

\subsection{KL-regularized PPO}

PPO \cite{schulman2017} receives a terminal reward-model score and token-level
shaping against the
frozen SFT policy:
\begin{align}
r_t^{\mathrm{KL}}
&=-\beta\left[\log\pi_\theta(y_t\mid s_t)
-\log\pi_{\mathrm{ref}}(y_t\mid s_t)\right],\\
R(x,y)&=r_\phi(x,y)+\sum_t r_t^{\mathrm{KL}} .
\end{align}
With generalized advantages $\hat A_t$ and
$\rho_t=\pi_\theta(y_t\mid s_t)/\pi_{\mathrm{old}}(y_t\mid s_t)$, the policy
loss is
\begin{equation}
\mathcal{L}_{\mathrm{PPO}}
=-\mathbb{E}_t\!\left[
\min\{\rho_t\hat A_t,
\operatorname{clip}(\rho_t,1-\epsilon,1+\epsilon)\hat A_t\}
\right].
\label{eq:ppo}
\end{equation}
The final path follows high-impact N+ implementation details
\cite{huang2024}: zero policy dropout, behavior log probabilities matched to
sampling temperature, a frozen SFT reference, an SFT-initialized reward model,
controlled scalar-head initialization, reward centering, an RM-initialized
critic, masked reward/advantage whitening, Adam $\epsilon=10^{-5}$, and
fixed-length sampling followed by EOS truncation. Responses without EOS
receive reward $-1$.

\begin{table*}[t]
\centering
\small
\caption{Final training profile. ``Token accuracy'' is teacher-forced
next-token accuracy over preferred completion tokens, not task accuracy.}
\label{tab:training}
\begin{tabularx}{0.96\textwidth}{@{}lYYYY@{}}
\toprule
\textbf{Stage} & \textbf{Data and budget} & \textbf{Adaptation} &
\textbf{Optimization} & \textbf{Outcome}\\
\midrule
SFT &
36,264 train; 1,917 val.; 4,096 total tokens &
LoRA $r=16$, $\alpha=32$, dropout 0 &
1 epoch; effective batch 32; LR $5{\times}10^{-6}$; cosine, 3\% warmup &
Train loss 1.0556; val.\ loss 1.1127; token accuracy 72.02\%\\
Reward model &
36,264 pairs; 1,917 val.\ pairs; 4,096 tokens &
SFT backbone; scalar head; LoRA $r=32$, $\alpha=64$ &
2 epochs; effective batch 64; LR $5{\times}10^{-6}$; centering coefficient
0.01 &
Val.\ loss 0.6166; pairwise accuracy 65.62\%; mean margin 0.3578\\
PPO &
24,018 unique prompts; 768-token rollouts &
SFT policy/reference; RM-initialized value model &
64 episodes/update; 4 PPO epochs; LR $3{\times}10^{-6}$; KL coefficient 0.07 &
6,400 episodes; 100 updates; objective KL 2.1648 at final update\\
\bottomrule
\end{tabularx}
\end{table*}

\section{Training Results}

\subsection{SFT assessment}

The one-epoch SFT run is stable. Validation loss changes only from 1.1183 at
step 250 to 1.1127 at step 1,134, while completion-token accuracy remains near
72\%. There is no late validation reversal suggesting material overfitting.
The sequence budget retains nearly all responses, and the effective batch,
learning rate, LoRA coverage, precision, and zero-dropout setting are
conservative for continued adaptation of a 0.5B model.

These metrics do not prove better open-ended responses, but they establish the
purpose required here: a stable HelpSteer-adapted policy and common reference.
Accordingly, this checkpoint is frozen for subsequent experiments rather than
rerun merely to obtain a lower training loss.

\subsection{Reward model}

The two-epoch reward model reaches 65.62\% overall pairwise accuracy. Accuracy
increases with preference strength: 61.23\%, 66.33\%, and 71.10\% for strengths
1, 2, and 3. Domain accuracy is uneven: 70.23\% on code, 62.91\% on general,
59.26\% on STEM, and 71.82\% on multilingual examples. This is usable as a
noisy preference signal, but not sufficiently reliable to serve as ground
truth, especially for correctness-sensitive STEM outputs.

An earlier custom reward model obtained 71.62\% validation accuracy yet
participated in a worse downstream PPO run. This is not evidence that lower
accuracy is desirable; it shows that aggregate in-distribution rank accuracy
does not capture reward scale, critic initialization, EOS handling, policy
optimization, or behavior on policy-generated text.

\subsection{PPO dynamics}

The selected PPO run was configured for 12,000 episodes and stopped for full
evaluation after 6,400 episodes and 100 optimizer updates. Average
response-level objective KL was 1.8278 and the final value was 2.1648.
This statistic aggregates over a response and is not directly comparable to
per-token KL targets. Old-to-new update KL remained much smaller
(mean 0.000640), the mean clip fraction was 0.00557, value loss fell from
25.32 to 5.50, and no empty-response collapse occurred. Optimization was
numerically stable, but stability alone did not prevent reward exploitation.

\section{Evaluation}

\subsection{Protocol}

Base, SFT, and PPO each generate one response for all 2,017 validation prompts
using the same tokenizer, prompt formatting, 3,072-token prompt limit,
1,024-token response cap, temperature 0.7, and top-$p=1$. The same epoch-two
reward model scores all prompt-response pairs. Pairwise tables derive from one
shared response table, preventing comparisons from using different samples.
Because the PPO training judge is also the evaluator, learned reward is treated
as a diagnostic rather than independent human preference.

\begin{table*}[t]
\centering
\small
\caption{Full policy-suite evaluation on 2,017 prompts.}
\label{tab:evaluation}
\begin{tabular}{@{}lrrrrrr@{}}
\toprule
\textbf{Policy} & \textbf{Three-way wins} & \textbf{Win rate} &
\textbf{Mean reward} & \textbf{Mean tokens} & \textbf{Median tokens} &
\textbf{Cap hit}\\
\midrule
Base & 718 & 35.60\% & 0.0803 & 398.0 & 331 & 8.82\%\\
SFT & 508 & 25.19\% & 0.0652 & 436.5 & 371 & 13.39\%\\
PPO & 775 & 38.42\% & 0.7300 & 577.8 & 520 & 27.42\%\\
Tie & 16 & 0.79\% & -- & -- & -- & --\\
\bottomrule
\end{tabular}
\end{table*}

\begin{table*}[t]
\centering
\small
\caption{Pairwise learned-reward results and PPO--Base domain breakdown. Win
rates include ties in the denominator.}
\label{tab:pairwise}
\begin{tabular}{@{}lrrrrr@{}}
\toprule
\textbf{Comparison} & \textbf{Left wins} & \textbf{Right wins} &
\textbf{Ties} & \textbf{Right win rate} & \textbf{Mean right--left delta}\\
\midrule
Base vs SFT & 1,158 & 837 & 22 & 41.50\% & $-0.0151$\\
Base vs PPO & 981 & 1,027 & 9 & 50.92\% & $+0.6497$\\
SFT vs PPO & 840 & 1,164 & 13 & 57.71\% & $+0.6648$\\
\midrule
\textbf{PPO vs Base by domain} & \textbf{Base wins} & \textbf{PPO wins} &
\textbf{Ties} & \textbf{PPO win rate} & \textbf{Prompts}\\
Code & 250 & 188 & 0 & 42.92\% & 438\\
General & 399 & 529 & 3 & 56.82\% & 931\\
STEM & 126 & 118 & 1 & 48.16\% & 245\\
Multilingual & 206 & 192 & 5 & 47.64\% & 403\\
\bottomrule
\end{tabular}
\end{table*}

PPO narrowly crosses 50\% against Base under the learned reward model, with the
aggregate advantage concentrated in general prompts. It remains below Base on
code and slightly below on STEM and multilingual prompts. The median
PPO-minus-Base reward delta is only $+0.0254$, whereas the mean is $+0.6497$;
a positive tail, including some flawed high-reward responses, raises the mean.

\subsection{Stopping, repetition, and qualitative audit}

PPO is substantially longer and more likely to exhaust the generation budget.
The fraction of responses whose repeated word-level 4-gram mass exceeds 25\%
is 10.11\% for Base, 20.08\% for SFT, and 31.88\% for PPO. At the stricter 50\%
threshold, the rates are 3.02\%, 8.48\%, and 15.82\%. Thus, the learned-reward
gain comes with an observable degradation in stopping and repetition.

Every evaluation row is assigned a deterministic triage label from reward
margins, EOS/cap status, and repeated 4-grams. The rules identify 8 likely
genuine PPO wins, 354 modest clean wins, 924 cases needing manual review, 228
repetition risks, 151 reward-model false-positive risks, 288 severe repetition
failures, and 64 strong PPO regressions. These labels prioritize review; they
are neither human judgments nor LLM-as-judge verdicts.

The public artifact contains all 2,017 full prompts and responses, with
secret-shaped strings redacted. A balanced 100-example subset contains 50
positive and 50 negative cases across all domains. Inspection finds real local
gains, such as better multi-part coverage and compact recovery from meandering
answers, alongside repeated loops, irrelevant continuation, malformed code,
fabricated citations, and incorrect scientific procedures. Publishing both
sides prevents the proxy score from becoming the entire account of the model.

\section{Discussion}

The strongest supported result is a systems-and-evaluation result. The project
connects long-context preprocessing, parameter-efficient SFT, pairwise reward
modeling, online PPO, checkpoint-safe execution, full-set generation, and
transparent audit on a diverse preference dataset. PPO changes behavior and
produces useful local gains, but the same run demonstrates how a numerically
stable policy can exploit a limited scalar judge.

The central limitations are:
\begin{enumerate}
  \item the reward model is both training signal and final quantitative judge;
  \item no blinded, multi-rater human preference study is reported;
  \item one stochastic generation per policy does not measure decoding
  variance;
  \item response length confounds reward, completeness, and repetition;
  \item a 0.5B policy and reward model have limited factual and long-context
  capacity; and
  \item the selected PPO checkpoint was practical rather than chosen by an
  independent human metric.
\end{enumerate}

\subsection{Controlled DPO extension}

The next experiment replaces reward-model optimization with Direct Preference
Optimization (DPO) \cite{rafailov2023} while keeping the SFT checkpoint,
preference records, tokenizer, 4,096-token budget, LoRA targets, and evaluation
prompts fixed:
\begin{align}
\Delta_\theta
&=\log\frac{\pi_\theta(y^+\mid x)}{\pi_{\mathrm{ref}}(y^+\mid x)}
-\log\frac{\pi_\theta(y^-\mid x)}{\pi_{\mathrm{ref}}(y^-\mid x)},\\
\mathcal{L}_{\mathrm{DPO}}
&=-\mathbb{E}\log\sigma\!\left(\beta\Delta_\theta\right).
\end{align}
The primary design uses standard sigmoid DPO with $\beta=0.1$, one epoch,
effective batch 32, LoRA rank 16, zero dropout, and precomputed reference log
probabilities. Each non-tied pair appears once. HelpSteer3 preference strength
is retained for stratified evaluation but is not silently interpreted as a
linear utility weight; an explicit strength-weighted sensitivity run can be
reported separately. The comparison will test whether direct offline
preference learning avoids the reward hacking and rollout cost observed under
PPO at this scale.

Further work should add blinded human review, repeated sampled generations and
paired confidence intervals, reward-model hard negatives drawn from observed
PPO failures, stronger correctness checks for code and STEM, and multi-metric
checkpoint selection. Scaling the policy is justified only after the evaluator
is less exploitable.

\section{Conclusion}

One-epoch continued SFT provides a stable HelpSteer-adapted reference, the
reward model learns a meaningful but imperfect ranking signal, and PPO
produces the strongest learned-reward result in the project. Yet the same
evaluation reveals longer, frequently capped, and substantially more
repetitive responses. The appropriate conclusion is neither that PPO failed
nor that it solved alignment. The project shows that a small-model RLHF stack
can be built and investigated rigorously, and that credible reporting requires
training metrics, full responses, and failure analysis together. DPO now
provides the cleanest controlled next comparison.

\balance
\begin{thebibliography}{99}

\bibitem{stiennon2020}
N. Stiennon et al.
\newblock Learning to summarize with human feedback.
\newblock \emph{NeurIPS}, 2020.
\newblock \url{https://arxiv.org/abs/2009.01325}.

\bibitem{ouyang2022}
L. Ouyang et al.
\newblock Training language models to follow instructions with human feedback.
\newblock \emph{NeurIPS}, 2022.
\newblock \url{https://arxiv.org/abs/2203.02155}.

\bibitem{huang2024}
S. Huang, M. Noukhovitch, A. Hosseini, K. Rasul, W. Wang, and L. Tunstall.
\newblock The N+ Implementation Details of RLHF with PPO.
\newblock arXiv:2403.17031, 2024.
\newblock \url{https://arxiv.org/abs/2403.17031}.

\bibitem{wang2025}
Z. Wang et al.
\newblock HelpSteer3-Preference: Open Human-Annotated Preference Data across
Diverse Tasks and Languages.
\newblock arXiv:2505.11475, 2025.
\newblock \url{https://arxiv.org/abs/2505.11475}.

\bibitem{hu2022}
E. J. Hu et al.
\newblock LoRA: Low-Rank Adaptation of Large Language Models.
\newblock \emph{ICLR}, 2022.
\newblock \url{https://arxiv.org/abs/2106.09685}.

\bibitem{schulman2017}
J. Schulman, F. Wolski, P. Dhariwal, A. Radford, and O. Klimov.
\newblock Proximal Policy Optimization Algorithms.
\newblock arXiv:1707.06347, 2017.
\newblock \url{https://arxiv.org/abs/1707.06347}.

\bibitem{qwen25}
Qwen Team.
\newblock Qwen2.5 Technical Report.
\newblock arXiv:2412.15115, 2024.
\newblock \url{https://arxiv.org/abs/2412.15115}.

\bibitem{rafailov2023}
R. Rafailov, A. Sharma, E. Mitchell, C. D. Manning, S. Ermon, and C. Finn.
\newblock Direct Preference Optimization: Your Language Model is Secretly a
Reward Model.
\newblock \emph{NeurIPS}, 2023.
\newblock \url{https://arxiv.org/abs/2305.18290}.

\bibitem{trl}
L. von Werra et al.
\newblock TRL: Transformer Reinforcement Learning.
\newblock Hugging Face, 2020--2026.
\newblock \url{https://huggingface.co/docs/trl/}.

\end{thebibliography}

\end{document}
